% meridianos.tex

\chapter{Meridianos}

Para representar uno o más meridianos en la esfera, tenemos que utilizar el comando:
\begin{center}
  {\small\texttt{\textbackslash dibujaTikZEsfera[true o false]}}
\end{center}
Para representar los meridianos, el comando Lua\LaTeX\ anterior necesita saber si la
esfera es transparente o no. Esto es necesario, pues en caso de que lo sea, se pueden
ver los meridianos que serían tapados por la esfera.

Por tanto, hay que indicar, entre corchetes, \texttt{true} si es transparente o
\texttt{false}, si es opaca.
Si no se detallara ningún argumento, el valor por defecto sería \texttt{false} lo que
indicaría que la esfera es opaca, por lo que no dejaría ver los elementos que se
encuentran tapados.

Con el comando anterior se leen, además de \texttt{esfera}, oros campos o subtablas,
como \texttt{'observador'}, \texttt{'meridprops'} y \texttt{'meridianos'} para representar
meridianos, pues su representación dependería del ángulo de observación de la esfera:

\begin{enumerate}
\item El campo \texttt{'observador'} define la vista de un observador situado muy lejos
  de la esfera, definida por las coordenadas angulares esféricas polar $\theta$ y
  azimutal $\phi$ en grados:
  \begin{itemize}
    \setlist[itemize]{noitemsep}\vspace*{-1ex}
  \item \texttt{thetaD}: Ángulo polar de observación de la esfera, en grados.
  \item \texttt{phiD}: Ángulo azimutal de observación, en grados.
  \end{itemize}
\item El campo \texttt{'meridprops'} define propiedades visuales por defecto del dibujo
  de los meridianos, y otras de la parte \emph{visible} y la \emph{invisible} de los
  meridianos (si la esfera fuera opaca):
  \begin{itemize}
    \setlist[itemize]{noitemsep}\vspace*{-1ex}
  \item \texttt{loops}: Cuántas veces se repite el bucle de trazado de cada meridano
    (de cuantos puntos consta cada meridiano).
  \item \texttt{color\_vis}: Color reservado para el dibujo de la parte \emph{visible}.
  \item \texttt{lw\_vis}: Grosor de línea de la parte \emph{visible}.
  \item \texttt{color\_novis}: Color reservado para el dibujo de la parte
    \emph{invisible}.
  \item \texttt{lw\_novis}: Grosor de línea de la parte \emph{invisible} (dimensión
    \LaTeX , preferiblemente en puntos (\unit{pt}))
  \end{itemize}  
\item El campo \texttt{'meridianos'}, define el ángulo azimutal $\phi$ que define la
  longitud en grados del meridiano:
  \begin{itemize}
        \setlist[itemize]{noitemsep}\vspace*{-1ex}
        \item \texttt{phiD}: Ángulo azimutal que define el meridiano, en grados (número)
        \end{itemize}
        Además del valor anterior, se puede destacar más o menos uno o más meridianos
        con respecto a los demás, añadiendo todos o alguno de los valores del
        campo \texttt{'meridprops'}. Ya veremos esto en algún ejemplo.
\end{enumerate}

Veamos uno por uno los campos en los ficheros que mostraremos en los siguientes
ejemplos, \texttt{esf\_merid\_01.lua} y \texttt{esf\_merid\_02.lua}:

\vspace{1ex}
CAMPO \texttt{'observador'}:

\begin{tabular}{|@{\hspace{-0.8em}}c|}\hline
  \vphantom{\rule{0.4pt}{3ex}}
  \begin{lstlisting}[style=myLuastyle]
    observador = {thetaD = 65, phiD = 15},
  \end{lstlisting}
  \\[1ex]
  \hline
\end{tabular}

\vspace{1ex}
CAMPO \texttt{'meridprops'}:

\begin{tabular}{|@{\hspace{-0.8em}}c|}\hline
  \vphantom{\rule{0.4pt}{3ex}}
  \begin{lstlisting}[style=myLuastyle]
   meridprops = {
      loops = 120, color_vis = "black!60", lw_vis="0.4pt",
      color_novis = "black!30", lw_novis="0.3pt",
   },
  \end{lstlisting}
  \\[1ex]
  \hline
\end{tabular}

\vspace{1ex}
CAMPO \texttt{'meridianos'}:

\begin{tabular}{|@{\hspace{-0.8em}}c|}\hline
  \vphantom{\rule{0.4pt}{3ex}}
  \begin{lstlisting}[style=myLuastyle]
   meridianos = {
     {phiD = 0,}, {phiD = 30}, {phiD = 60},
     {phiD = 90}, {phiD = 120}, {phiD = 150},
     {phiD = 180}, {phiD = -30}, {phiD = -60},
     {phiD = -90}, {phiD = -120},{phiD = -150}
   },
  \end{lstlisting}
  \\[1ex]
  \hline
\end{tabular}


En el fichero \texttt{esf\_merid\_01.lua} se define una esfera plana:

\vspace{1ex}
\begin{tabular}{|@{\hspace{-1em}}c|}\hline
  \vphantom{\rule{0.4pt}{3ex}}
  \begin{lstlisting}[basicstyle=\ttfamily\footnotesize, language=TeX]
   esfera = {radio=2.0, draw="black!25", lw="0.6pt", fill="black!8", opac=1.0,},
  \end{lstlisting}
  \\[1ex]
  \hline
\end{tabular}

La rutina de trazado de meridianos se encarga de dibujar primero la parte
\emph{invisible} (en caso de que la esfera sea translúcida), y solo después,
la parte \emph{visible}.

Para no repetir el mismo código, resaltaremos el fichero de datos y el comando
\LaTeX\ que utilizamos Para una esfera opaca sería:

\vspace{1ex}
\begin{tabular}{|@{\hspace{-1em}}c|}\hline
  \vphantom{\rule{0.4pt}{3ex}}
  \begin{lstlisting}[basicstyle=\ttfamily\footnotesize, language=TeX]
    \cargaDatos{datoslua/esf_merid_01.lua}
        ...   
        \dibujaTikzEsfera
        ...
  \end{lstlisting}
  \\[1ex]
  \hline
\end{tabular}

El fichero \texttt{esf\_merid\_02.lua}, define una esfera sombreada (el resto de
campos es el mismo que en el fichero \texttt{esf\_merid\_01.lua}:

\vspace{1ex}
\begin{tabular}{|@{\hspace{-1em}}c|}\hline
  \vphantom{\rule{0.4pt}{3ex}}
  \begin{lstlisting}[basicstyle=\ttfamily\footnotesize, language=TeX]
   esfera = {radio=2.0, draw="black!25", lw="0.6pt", fill="black!8", opac=1.0,},
   smbresfera = {color="gray!1", opac=0.2},
  \end{lstlisting}
  \\[1ex]
  \hline
\end{tabular}

A continuación se muestran las dos esferas:

 \begin{figure}[ht]
  \centering
  \begin{minipage}{0.45\textwidth}
    \cargaDatos{datoslua/esf_merid_01.lua}
    \begin{tikzpicture}[scale=1.0,]
      \dibujaTikzEsfera
    \end{tikzpicture}
    \caption{Esfera plana.}
  \end{minipage}
  \hfill
  \begin{minipage}{0.45\linewidth}
    \cargaDatos{datoslua/esf_merid_02.lua}
     \centering
     \begin{tikzpicture}[scale=1.0,]
       \dibujaTikzEsfera
     \end{tikzpicture}
     \caption{Esfera sombreada.}
   \end{minipage}
   \caption{Esferas opacas.}
 \end{figure}


Ahora dibujamos las mismas esferas y meridianos, pero considerando que son translúcidas
(añadiendo \texttt{true} al comando: \texttt{\textbackslash dibujaTikzEsfera[true]}

\begin{figure}[ht]
  \centering
  \begin{minipage}[b]{0.45\textwidth}
    \cargaDatos{datoslua/esf_merid_01.lua}
    \begin{tikzpicture}[scale=1.0,]
      \dibujaTikzEsfera[true]
    \end{tikzpicture}
    \caption{Esfera plana.}
  \end{minipage}
  \hfill
  \begin{minipage}[b]{0.45\linewidth}
    \cargaDatos{datoslua/esf_merid_02.lua}
    \centering
    \begin{tikzpicture}[scale=1.0,]
      \dibujaTikzEsfera[true]
    \end{tikzpicture}
    \caption{Esfera sombreada.}
  \end{minipage}
  \caption{Esferas translúcidas.}
\end{figure}

Otras figuras:

En la siguientes figuras opacas resaltamos un meridiano.
Para ello, utilizamos el fichero \texttt{esf\_meridres\_01.lua} para la esfera
gris:

\vspace{1ex}
\begin{tabular}{|@{\hspace{-1em}}c|}\hline
  \vphantom{\rule{0.4pt}{3ex}}
  \begin{lstlisting}[basicstyle=\ttfamily\footnotesize, language=TeX]
    meridprops = {
      loops = 120, color_vis = "black!60", lw_vis="0.4pt",
      color_novis = "black!30", lw_novis="0.3pt",
    },
  \end{lstlisting}
  \\[1ex]
  \hline
\end{tabular}

Se puede observar que para esta esfera se resaltan de rojo los puntos visibles y
naranja los invisibles, los meridianos \qty{0}{\degree} y \qty{180}{\degree}:

\vspace{1ex}
\begin{tabular}{|@{\hspace{-1em}}c|}\hline
  \vphantom{\rule{0.4pt}{3ex}}
  \begin{lstlisting}[basicstyle=\ttfamily\footnotesize, language=TeX]
    meridianos = {
      {phiD = 0,
        color_vis="red", lw_vis="0.8pt", color_novis="orange", lw_novis="0.8pt"},
      {phiD = 30}, {phiD = 60}, {phiD = 90}, {phiD = 120}, {phiD = 150},
      {phiD = 180,
        color_vis="red", lw_vis="0.8pt", color_novis="orange", lw_novis="0.8pt"},
      {phiD = -30}, {phiD = -60}, {phiD = -90},  {phiD = -120}, {phiD = -150}
    },
  \end{lstlisting}
  \\[1ex]
  \hline
\end{tabular}

Para la verde utilizamos \texttt{esf\_meridres\_02.lua}:

\vspace{1ex}
\begin{tabular}{|@{\hspace{-1em}}c|}\hline
  \vphantom{\rule{0.4pt}{3ex}}
  \begin{lstlisting}[basicstyle=\ttfamily\footnotesize, language=TeX]
    meridprops = {
      loops = 120, color_vis = "black!80", lw_vis="0.6pt",
      color_novis = "black!40", lw_novis="0.5pt",
    },
  \end{lstlisting}
  \\[1ex]
  \hline
\end{tabular}

Para la esfera verde se resaltan de rojo los puntos visibles y en  naranja los
invisibles, los meridianos \qty{-30}{\degree} y \qty{150}{\degree}:

\vspace{1ex}
\begin{tabular}{|@{\hspace{-1em}}c|}\hline
  \vphantom{\rule{0.4pt}{3ex}}
  \begin{lstlisting}[basicstyle=\ttfamily\footnotesize, language=TeX]
    meridianos = {
      {phiD = 0,}, {phiD = 30}, {phiD = 60}, {phiD = 90}, {phiD = 120},
      {phiD = 150,
        color_vis="red", lw_vis="1pt", color_novis="orange", lw_novis="1pt"},
      {phiD = 180},
      {phiD = -30,
        color_vis="red", lw_vis="1pt", color_novis="orange", lw_novis="1pt"},
      {phiD = -60}, {phiD = -90}, {phiD = -120}, {phiD = -150}
    },
  \end{lstlisting}
  \\[1ex]
  \hline
\end{tabular}

\begin{figure}[ht]
  \centering
  \begin{minipage}[b]{0.45\textwidth}
    \cargaDatos{datoslua/esf_meridres_01.lua}
    \begin{tikzpicture}[scale=1.0,]
      \dibujaTikzEsfera
    \end{tikzpicture}
    \caption{Esfera gris.}
  \end{minipage}
  \hfill
  \begin{minipage}[b]{0.45\linewidth}
    \cargaDatos{datoslua/esf_meridres_02.lua}
    \centering
    \begin{tikzpicture}[scale=1.0,]
      \dibujaTikzEsfera
    \end{tikzpicture}
    \caption{Esfera verde.}
  \end{minipage}
  \caption{Esferas opacas con dos meridianos resaltados.}
\end{figure}


\begin{figure}[ht]
  \centering
  \begin{minipage}[b]{0.45\textwidth}
    \cargaDatos{datoslua/esf_meridres_01.lua}
    \begin{tikzpicture}[scale=1.0,]
      \dibujaTikzEsfera[true]
    \end{tikzpicture}
    \caption{Esfera gris.}
  \end{minipage}
  \hfill
  \begin{minipage}[b]{0.45\linewidth}
    \cargaDatos{datoslua/esf_meridres_02.lua}
    \centering
    \begin{tikzpicture}[scale=1.0,]
      \dibujaTikzEsfera[true]
    \end{tikzpicture}
    \caption{Esfera verde.}
  \end{minipage}
  \caption{Esferas translúcidas con dos meridianos resaltados.}
\end{figure}












%%% Local Variables:
%%% coding: utf-8
%%% mode: latex
%%% TeX-engine: luatex
%%% TeX-master: "../elementos_esfera.tex"
%%% End:


